\chapter{Двухкомпонентная source-parent архитектура дополнительных щелей}
\label{ch:v8-extra-mass-two-source-parent}

\section{Наследованная следовая жёсткость}

Сохраним ортогональный бесследовый базис предыдущего гейта
\begin{equation}
 A=P_u-P_d,\qquad B=P_u+P_d-3P_Y,
 \qquad \Tr A^2=12,\quad \Tr B^2=48,\quad \Tr AB=0.
 \label{eq:v8-extra-mass-two-source-basis}
\end{equation}
Для $h^0=aA+bB$ минимальный parent, не вводящий новой матрицы жёсткости,
есть
\begin{equation}
 V_{j_A,j_B}(a,b)=\frac12\Tr(h^0)^2-j_Aa-j_Bb
 =6a^2+24b^2-j_Aa-j_Bb.
 \label{eq:v8-extra-mass-two-source-potential}
\end{equation}
Его гессиан совпадает с Gram-метрикой центральной плоскости:
\begin{equation}
 \Hess V=\begin{pmatrix}12&0\\0&48\end{pmatrix}>0.
 \label{eq:v8-extra-mass-two-source-hessian}
\end{equation}
Тем самым функционал строго выпуклый и имеет единственный глобальный минимум.
Квадрат завершается точно:
\begin{equation}
 V=6\left(a-\frac{j_A}{12}\right)^2
 +24\left(b-\frac{j_B}{48}\right)^2
 -\frac{j_A^2}{24}-\frac{j_B^2}{96}.
 \label{eq:v8-extra-mass-two-source-completed-square}
\end{equation}
Норм-квадраты отображений момента дают стандартный источник положительных
квадратичных функционалов на пространствах представлений колчанов
\cite{v8-two-source-harada-wilkin}; spectral action также накладывает
ограничения на скалярный потенциал лишь после задания конечной геометрии
\cite{v8-two-source-krajewski}. Эти результаты обосновывают архитектуру, но
не порождают $j_A,j_B$ в текущем родителе.

\section{Биекция между источниками и щелями}

Стационарная точка равна
\begin{equation}
 a_*=\frac{j_A}{12},\qquad b_*=\frac{j_B}{48}.
 \label{eq:v8-extra-mass-two-source-minimum}
\end{equation}
Подстановка в карту щелей предыдущего гейта даёт
\begin{equation}
 \begin{pmatrix}\Delta_u\\\Delta_d\end{pmatrix}
 =\frac1{12}
 \begin{pmatrix}1&1\\-1&1\end{pmatrix}
 \begin{pmatrix}j_A\\j_B\end{pmatrix}.
 \label{eq:v8-extra-mass-two-source-gap-map}
\end{equation}
Определитель этой карты равен
\begin{equation}
 \det\frac{\partial(\Delta_u,\Delta_d)}
 {\partial(j_A,j_B)}=\frac1{72}\ne0,
 \label{eq:v8-extra-mass-two-source-determinant}
\end{equation}
поэтому обратная карта существует на всей плоскости:
\begin{equation}
 j_A=6(\Delta_u-\Delta_d),\qquad
 j_B=6(\Delta_u+\Delta_d).
 \label{eq:v8-extra-mass-two-source-inverse}
\end{equation}
В частности,
\begin{equation}
 (\Delta_u,\Delta_d)=(\delta,\delta)
 \Longleftrightarrow (j_A,j_B)=(0,12\delta),
 \label{eq:v8-extra-mass-two-source-equal-witness}
\end{equation}
а асимметричный свидетель $(1,2)$ требует $(-6,18)$.

\section{Минимальность и статус происхождения}

Один вещественный источник имеет образ размерности не более одного и не
может параметризовать открытую часть двумерного симплекса щелей. Два
источника достаточны по ненулевому определителю
\eqref{eq:v8-extra-mass-two-source-determinant}; следовательно, их число
минимально. Общая положительная матрица жёсткости внесла бы три новых
вещественных коэффициента,
\begin{equation}
 K=\begin{pmatrix}k_{AA}&k_{AB}\\k_{AB}&k_{BB}\end{pmatrix}>0,
 \label{eq:v8-extra-mass-two-source-general-stiffness}
\end{equation}
но здесь $K=\operatorname{diag}(12,48)$ уже фиксирована незавешенным следом.

Операторы $A,B$ центральны, вещественны и чётны, поэтому линейный член
калибровочно и grading-совместим. Однако совместимость не отвечает на вопрос,
какой существующий скалярный носитель создаёт две независимые компоненты.
Пара $(j_A,j_B)$ распадается на один модуль и одно проективное направление;
оба пока являются новыми данными. Физические KMS-веса дополнительно зависят
от $\beta\Delta_u$ и $\beta\Delta_d$, так что архитектура не фиксирует
температурный или временной масштаб.

Итоговые реестры:
\begin{equation}
 N_{\rm architecture}=9/9,
 \qquad N_{\rm source\ origin}=0/2.
 \label{eq:v8-extra-mass-two-source-ledgers}
\end{equation}

\begin{theorem}[Минимальная двухисточниковая реализация]
При наследованной следовой жёсткости $\operatorname{diag}(12,48)$ функционал
$V_{j_A,j_B}$ строго выпуклый. Его единственный минимум задаёт биекцию между
двумя вещественными source-компонентами и произвольной парой центральных
щелей. Двух источников достаточно и необходимо.
\end{theorem}

\begin{nogocriterion}[Форма не равна происхождению]
Нельзя считать пару щелей выведенной только потому, что положительная
следовая метрика делает source-parent устойчивым. До предъявления двух
типизированных скалярных носителей либо единого носителя с двумерным образом
$j_A,j_B$ остаются независимыми входными данными.
\end{nogocriterion}

\begin{protocol}\begin{enumerate}
\item жёсткость: \textbf{$\operatorname{diag}(12,48)$};
\item новых stiffness-параметров: \textbf{ноль};
\item минимум: \textbf{$(j_A/12,j_B/48)$};
\item determinant source-to-gap: \textbf{$1/72$};
\item произвольная пара щелей реализуема: \textbf{да};
\item один источник достаточен: \textbf{нет};
\item два источника достаточны: \textbf{да};
\item gauge/grading compatibility: \textbf{да};
\item architecture-ledger: \textbf{$9/9$};
\item source-origin: \textbf{$0/2$};
\item следующий гейт: \textbf{классификация существующих скалярных носителей двух источников}.
\end{enumerate}\end{protocol}

Машинный сертификат сохранён в
\path{s2t_v8_baryon_c0_singlet_triplet_central_gap_isotypic_channel_extra_edge_mass_two_source_parent_architecture_gate_results.json}.

\begin{thebibliography}{9}
\bibitem{v8-two-source-harada-wilkin} M.~Harada, G.~Wilkin,
\emph{Morse Theory of the Moment Map for Representations of Quivers},
arXiv:0807.4734.
\bibitem{v8-two-source-krajewski} T.~Krajewski,
\emph{Constraints on the Scalar Potential from Spectral Action Principle},
arXiv:hep-th/9803199.
\end{thebibliography}